\section{Introduction}
\label{sec:introduction}

Organic semiconductors provide a route to electronic devices that differ in
important ways from conventional inorganic semiconductor technology. They are
soft molecular materials, can often be processed under comparatively mild
conditions, and are suitable for deposition over large areas and on lightweight
or mechanically flexible substrates. These features make them attractive for
future display technologies, including foldable mobile devices and large-area
screens that can be bent or rolled. At the same time, their electronic
properties can be influenced by molecular design, film growth, interfaces, and
processing conditions. Organic semiconductors are therefore relevant not only
for flexible displays, but also for sensors, low-cost integrated circuits, and
optoelectronic devices \cite{Bruetting2005,Sirringhaus2005DevicePhysics}.

This flexibility in materials and processing is also one of the reasons why
organic electronic devices are difficult to characterize quantitatively. Charge
transport is affected by processes on several length scales: molecular packing
and localized states determine the local transport landscape, while grain
boundaries, interfaces, contacts, and device geometry shape the macroscopic
current. Charge-carrier mobility is commonly used to compare materials and
devices, but it is not a quantity that is measured directly. It is extracted
from electrical data with the help of a transport model. The resulting value
therefore depends on the assumptions of that model, on the voltage range used
for the extraction, and on nonidealities of the device itself.

Organic field-effect transistors are widely used for such mobility
measurements. In a conventional two-terminal measurement, however, the measured
voltage includes both the voltage drop across the semiconductor channel and the
voltage losses at the metal--semiconductor contacts. The gradual-channel model
usually applied to transfer curves assumes that these contact contributions are
small compared with the channel resistance. This assumption can be problematic
for organic thin films. Charge injection at the metal--organic interface depends
on energetic alignment, local morphology, trap density, and gate bias
\cite{Natali2012ChargeInjection}. As a consequence, the slope of a measured
transfer characteristic may reflect changes in both the channel and the
contacts. If the ideal transistor model is then applied without further checks,
the extracted mobility can describe an apparent device response rather than the
transport in the semiconductor film. In severe cases, conventional
current--voltage analysis has been shown to overestimate mobility by one order
of magnitude or more \cite{Bittle2016MobilityOverestimation}.

The gated van der Pauw method is one way to reduce this ambiguity. It combines
field-effect modulation with a four-terminal sheet-conductance measurement. A
current is driven through one pair of contacts, while a separate pair of
nominally non-injecting contacts senses the voltage. The voltage drop at the
current-injecting contacts is therefore not included directly in the measured
van der Pauw resistance. By sweeping the gate voltage, the carrier density in
the accumulation layer is changed, and the corresponding sheet conductance can
be used to determine a threshold voltage and an effective thin-film mobility
that is less sensitive to contact voltage losses.

The word contact-independent must nevertheless be used with care. The current
still has to enter and leave the organic semiconductor through metallic
contacts, and poor injection can limit the accessible measurement range or
require large driving voltages. The method only removes the contact voltage
drops from the resistance used for the sheet-conductance extraction, provided
that the imposed current is stable and the voltage probes draw negligible
current. The extracted mobility also remains a macroscopic device parameter. It
can include the influence of energetic disorder, traps, grain boundaries, and
spatial inhomogeneities, and should not automatically be identified with an
intrinsic molecular or intra-grain mobility.

The aim of this thesis is to characterize organic thin-film field-effect
devices in a way that separates, as far as possible, the accumulated
semiconductor sheet from contact-dominated device response. For this purpose,
conventional current--voltage measurements are combined with a current-driven
gated van der Pauw method. The work begins with the characterization of a
prototype van der Pauw measurement setup and continues with the fabrication and
analysis of molecular semiconductor devices prepared using physical vapor
deposition, chemical vapor deposition, and spin coating. Special attention is
paid to the local gate overdrive, to the dependence of the extracted quantities
on the imposed measurement current, to residual series resistance in the full
current path, and to transient effects that can influence the measured response.
In the experimental chapters, these checks are applied to C8-BTBT OFETs in which
two-terminal data confirm field-effect operation but also reveal nonideal
features, whereas the gated van der Pauw analysis gives a current-independent
sheet conductance and a mobility of about
\SI{3.1}{\centi\metre\squared\per\volt\per\second}. The measurements are
therefore treated not as isolated electrical curves, but as controlled tests of
the conditions under which reliable transport parameters can be obtained.

